\banner{Design project 14: a drainage system for an irrigated clay field}\label{sec:proj14}
\textbf{Brief.} A 40-ha field of cotton ($EC_e$ threshold 7.7 dS\,m$^{-1}$) is irrigated with water of 1.5 dS\,m$^{-1}$ and has a barrier at 6 m. The design discharge is 7 mm\,d$^{-1}$ and the water table midway between drains is allowed to stand $h=0.5$ m above drain level; $K_1=0.8$, $K_2=0.3$ m\,d$^{-1}$. The designer chooses a drain depth, and then the spacing and the total pipe.

\begin{table}[!htbp]\centering\small\caption{Hooghoudt spacings for $q=7$ mm/d, $h=0.5$ m, $K_1=0.8$, $K_2=0.3$ m/d, barrier at 6 m.}
\begin{tabular}{rrrrr}\toprule
\hd{Drain depth (m)} & \hd{$d$ (m)} & \hd{$d_e$ (m)} & \hd{Spacing $L$ (m)} & \hd{Drain length (m/ha)}\\\midrule
1.2 & 4.8 & 1.9 & 21 & 472\\
1.5 & 4.5 & 1.9 & 21 & 475\\
1.8 & 4.2 & 1.9 & 21 & 478\\
2.1 & 3.9 & 1.8 & 21 & 482\\\bottomrule\end{tabular}\end{table}

\textbf{Reading.} Deeper drains lower the barrier-limited flow path only slightly, because the equivalent depth saturates as the drain nears the barrier, so the spacing rises little with depth; the extra excavation cost must be set against the small gain in spacing. At 1.8 m, $L=21$ m, so the field of 40 ha needs about $19$ km of pipe. The leaching requirement for cotton at an $EC_e$ of 7.7 is $0.041$, small, so the drainage is set by the water table rather than by salt.
\fig{g14_wt}{0.7}{Steady water table midway between two drains of spacing $L=21$ m (schematic elliptical shape, $h=0.5$ m).}
